% ============================================================
% preamble.tex —— 《烧掉数学书之后》全书公共导言
% 编译方式：xelatex（需安装 ctex 宏包，TeX Live / MiKTeX 均可）
% ============================================================
\usepackage{geometry}
\geometry{a4paper,left=2.6cm,right=2.6cm,top=2.7cm,bottom=3.0cm}

\usepackage{amsmath,amssymb,amsthm,mathtools}
\usepackage{array,booktabs,longtable,multirow}
\usepackage{enumitem}
\usepackage{graphicx}
\usepackage[most]{tcolorbox}
\usepackage{tikz}
\usetikzlibrary{arrows.meta,calc,decorations.markings}
\usepackage[hidelinks]{hyperref}
\hypersetup{
  bookmarksnumbered=true,
  pdftitle={烧掉数学书之后：和初中生一起重造数学},
  pdfauthor={〔作者姓名〕},
}

% ---------- 主题色 ----------
\definecolor{shaohong}{RGB}{178,34,34}   % 烧红
\definecolor{cheng}{RGB}{224,122,32}     % 橙：开胃问题
\definecolor{shenlan}{RGB}{41,84,150}    % 蓝：发明时刻
\definecolor{molv}{RGB}{30,132,73}       % 绿：动手实验
\definecolor{zi}{RGB}{108,52,131}        % 紫：烧脑时刻
\definecolor{zong}{RGB}{121,85,61}       % 棕：诚实声明
\definecolor{jin}{RGB}{163,124,34}       % 金：发明词典
\definecolor{hui}{RGB}{100,100,100}

% ---------- 全书栏目盒子 ----------
% 开胃问题：每章开场的具体小问题
\newtcolorbox{kaiwei}{enhanced,breakable,
  colback=cheng!6!white,colframe=cheng!90!black,
  fonttitle=\bfseries,title={\large 开胃问题},
  boxrule=0.9pt,arc=2.5mm,left=3mm,right=3mm,top=1.5mm,bottom=2mm}

% 发明时刻：本章被迫发明了什么（土名 → 正式名）
\newtcolorbox{faming}{enhanced,breakable,
  colback=shenlan!6!white,colframe=shenlan!85!black,
  fonttitle=\bfseries,title={\large 发明时刻},
  boxrule=0.9pt,arc=2.5mm,left=3mm,right=3mm,top=1.5mm,bottom=2mm}

% 发明词典：土名与正式术语对照
\newtcolorbox{cidian}{enhanced,breakable,
  colback=jin!7!white,colframe=jin!90!black,
  fonttitle=\bfseries,title={\large 发明词典},
  boxrule=0.9pt,arc=2.5mm,left=3mm,right=3mm,top=1.5mm,bottom=2mm}

% 工具箱：现场自造的通用工具
\newtcolorbox{gongju}{enhanced,breakable,
  colback=hui!7!white,colframe=hui!75!black,
  fonttitle=\bfseries,title={\large 工具箱},
  boxrule=0.9pt,arc=2.5mm,left=3mm,right=3mm,top=1.5mm,bottom=2mm}

% 动手实验：分级练习 ★ / ★★ / ★★★
\newtcolorbox{dongshou}{enhanced,breakable,
  colback=molv!6!white,colframe=molv!85!black,
  fonttitle=\bfseries,title={\large 动手实验},
  boxrule=0.9pt,arc=2.5mm,left=3mm,right=3mm,top=1.5mm,bottom=2mm}

% 烧脑时刻：章末选做挑战
\newtcolorbox{shaonao}{enhanced,breakable,
  colback=zi!6!white,colframe=zi!85!black,
  fonttitle=\bfseries,title={\large 烧脑时刻},
  boxrule=0.9pt,arc=2.5mm,left=3mm,right=3mm,top=1.5mm,bottom=2mm}

% 诚实声明：本书在哪里偷了懒
\newtcolorbox{chengshi}{enhanced,breakable,
  colback=zong!6!white,colframe=zong!85!black,
  fonttitle=\bfseries,title={诚实声明},
  boxrule=0.7pt,arc=2mm,left=3mm,right=3mm,top=1.5mm,bottom=2mm}

% 数学史话
\newtcolorbox{lishi}{enhanced,breakable,
  colback=shaohong!5!white,colframe=shaohong!80!black,
  fonttitle=\bfseries,title={数学史话},
  boxrule=0.7pt,arc=2mm,left=3mm,right=3mm,top=1.5mm,bottom=2mm}

% 结论框：重点结论（文字自动换行，不溢出版心）
\newtcolorbox{jielun}{enhanced,breakable,
  colback=shaohong!5!white,colframe=shaohong!75!black,
  boxrule=0.9pt,arc=2.5mm,left=3mm,right=3mm,top=2mm,bottom=2mm,
  before skip=8pt,after skip=8pt}

% 练习条目快捷方式
\newcommand{\yisi}{\item[$\star$]}
\newcommand{\yier}{\item[$\star\star$]}
\newcommand{\yisan}{\item[$\bigstar\bigstar\bigstar$]}
\setlist[itemize,1]{leftmargin=2.2em,itemsep=0.3em}

% ---------- 定理类环境 ----------
\theoremstyle{plain}
\newtheorem{dingli}{定理}[chapter]
\newtheorem{yinli}[dingli]{引理}
\newtheorem{tuilun}[dingli]{推论}
\newtheorem{mingti}[dingli]{命题}
\theoremstyle{definition}
\newtheorem{dingyi}{定义}[chapter]
\newtheorem{caixiang}{猜想}[chapter]
\newtheorem{li}{例}[chapter]
\renewcommand{\proofname}{证明}

% 小工具
\newcommand{\Keyword}[1]{\textbf{#1}}
\newcommand{\shao}{\textcolor{shaohong}{\textbf{烧}}}
\providecommand{\heart}{\heartsuit} % 序章自定义运算符号
